\section{Experiments}
\label{sec:experiments}

\subsection{Dataset and benchmark}
We use COCO~2017~\cite{lin2014coco}: $118{,}287$ training images, $5{,}000$
validation images, and $40{,}670$ held-out \texttt{test2017} images. Fine-tuning
uses \texttt{train2017}; the \texttt{val2017} split drives best-checkpoint
selection during training. Because \texttt{val2017} is therefore ``seen'' by the
training procedure, evaluating on it would leak. Our benchmark is instead a
\emph{stratified $1{,}000$-image subset of \texttt{test2017}} (fixed seed
$42$), which neither fine-tuning nor checkpoint selection ever touches. All
models are scored on this identical benchmark.

\subsection{Models compared}
Table~\ref{tab:models} lists the four model variants spanning the three
paradigms we benchmark. Two are our own CNN (the pretrained encoder and our
COCO-fine-tuned model); the other two are pretrained third-party systems
wrapped behind a unified \texttt{colorize(gray)} interface so they receive
identical inputs.

\begin{table}[t]
\centering
\caption{The four model variants evaluated on the shared benchmark.}
\label{tab:models}
\begin{tabular}{@{}llc@{}}
\toprule
Model & Paradigm & Ours? \\
\midrule
Zhang16 Pretrained (ECCV)      & CNN             & reimpl. \\
Zhang16 Fine-tuned (COCO 2017) & CNN             & yes \\
Zhang17 (automatic mode)       & Interactive CNN & no \\
DeOldify (NoGAN)               & GAN             & no \\
\bottomrule
\end{tabular}
\end{table}

\subsection{Why diffusion is excluded from the comparison}
\label{sec:experiments:no-diffusion}
Conditional diffusion is the natural fourth paradigm and the literature offers
several strong recipes---image-to-image diffusion in the spirit of
Palette~\cite{saharia2022palette}, and conditioning a pretrained text-to-image
backbone via ControlNet~\cite{zhang2023controlnet} adapters. We attempted the
ControlNet recipe but could not complete it under conditions that would yield
a meaningful comparison number, for a single, well-defined reason: every
well-trained openly-released diffusion colorization checkpoint we could find
is anchored to \emph{Stable Diffusion 2.1} as its image backbone, and
Stability~AI has deprecated the entire SD~2.x line upstream. The relevant
Hugging Face repositories (the \texttt{stabilityai/stable-diffusion-2*}
namespace) return HTTP~401 even with a valid token, and no public mirror
provides licensed access.

We surveyed three substitution paths and rejected each:
\begin{itemize}
  \item \textbf{An SD~1.5 brightness-ControlNet} (\texttt{ioclab/control\_v1p\_sd15\_\allowbreak brightness}).
        The conditioning channel is brightness, not chrominance; on our benchmark
        it produced PSNR ${\approx}6$\,dB---worse than a constant-gray output.
  \item \textbf{An SD~InstructPix2Pix colorization fine-tune}. The model emitted
        NaN tensors mid-denoise and required ${\sim}9$\,min per image on
        consumer hardware---neither a useful number nor a usable pipeline.
  \item \textbf{A Flux-based pipeline}. At 12\,B parameters it does not fit on
        the 6\,GB GTX~1660~SUPER we used for the rest of the study, breaking
        the ``identical conditions'' requirement.
\end{itemize}
Substituting any of these into Table~\ref{tab:results} would have placed a
misleading row next to faithfully-evaluated systems. We therefore restrict the
comparison to the three paradigms above and treat the diffusion gap as a
finding in its own right: the open-source diffusion-colorization stack was, in
2026, single-point-of-failure dependent on a backbone that has now been
withdrawn. We revisit this in Sec.~\ref{sec:discussion}.

\subsection{Metrics and uncertainty}
We report three complementary metrics, all computed in RGB against the ground
truth: PSNR (pixel fidelity), SSIM (structural similarity), and
LPIPS~\cite{zhang2018lpips} (learned perceptual distance; lower is better).
Because a single mean hides sampling variability, every metric is reported as
\texttt{mean [95\% CI]}, where the confidence interval is a percentile bootstrap
over the $1{,}000$ benchmark images ($10{,}000$ resamples, seed~$42$).

\subsection{Training configuration}
Our model is initialized from the ECCV encoder weights, but its $233$-bin
classification head is trained from scratch---the official head has $313$ bins
and is dimensionally incompatible with our gamut (Section~\ref{sec:method}).
Table~\ref{tab:hparams} lists the hyperparameters. Mixed precision was disabled
after the fp16 path overflowed to NaN mid-epoch; training therefore runs in
fp32 at batch size $4$, which fits the $6$\,GB GTX~1660~SUPER. We trained for
eight epochs in two phases: five batch-capped epochs followed by three
full-data epochs resumed from the prior checkpoint. Fig.~\ref{fig:curves} shows
the loss and validation-PSNR trajectories.

\begin{table}[t]
\centering
\caption{Fine-tuning hyperparameters.}
\label{tab:hparams}
\begin{tabular}{@{}ll@{}}
\toprule
Setting & Value \\
\midrule
Optimizer        & Adam, $\beta=(0.9,0.999)$, weight decay $0$ \\
Learning rate    & $2\times10^{-4}$, StepLR (step $20$, $\gamma=0.5$) \\
Batch size       & $4$ (fp32; AMP disabled) \\
Gradient clip    & $1.0$ (max-norm) \\
Input resolution & $256\times256$ \\
Quantization     & $Q=233$ bins; soft target, $\sigma=5$, $5$-NN \\
Rebalancing      & $\lambda=0.5$ \\
Decoder          & annealed mean, $T=0.38$ \\
Epochs           & $8$ ($5$ capped $+\,3$ full-data) \\
Hardware         & GTX~1660~SUPER (6\,GB), PyTorch~2.5.1+cu121 \\
\bottomrule
\end{tabular}
\end{table}

\begin{figure}[t]
  \centering
  \includegraphics[width=\columnwidth]{training_curves.png}
  \caption{Fine-tuning on COCO~2017. Left: training and validation loss; two
  validation spikes from pathological batches are off-scale and annotated.
  Right: validation PSNR climbs to $23.6$\,dB by epoch~8.}
  \label{fig:curves}
\end{figure}
